\documentclass[10pt]{article}

\usepackage[utf8]{inputenc}

\usepackage[T1]{fontenc}

\usepackage{amsmath}

\usepackage{amsfonts}

\usepackage{amssymb}

\usepackage[version=4]{mhchem}

\usepackage{stmaryrd}

\usepackage{bbold}



\begin{document}

В н у т р ен в я я п л о а д ь. Если $M$ метризовано и $f: M \rightarrow \mathbb{R}^{n}$ - локально изометрич. отображение, то возникает вопрос о соотношении между $L(M, f)$ и мерой Хаусдорфа $H_{k}(M)$. Когда $M$ - двумерное жногообразие озраниченной кривизны, то $L(M, f)=$ $=H_{2}(M)$. Вообще же непрерывное $f: M \rightarrow \mathbb{R} n$ индуцирует на компонентах связности $f^{-1}(u), u \in \mathbb{R}^{n}$, обобщенную метрику $\rho$, отличающуюся возможностью $\rho(x, y)=\infty$. Конструкция $k$-мерной меры Хаусдорфа, примененная к $\rho$, дает характеристику, к-рую можно принять за внутреннюю П. погружения. Для липшицевых $f$ она совпадает с $L(M, f)(\mathrm{cm},[11])$.



Массы По токов и вар ио о ра а йй. Интегрирование $k$-форм по $k$-мерному кусочно гладко вложенному в $\mathbb{R}^{n}$ ориентированному многообразию $M$ приводит к потоку $T_{M \varphi}=\int_{M}\langle\varphi(x), v(x)\rangle d F(x)-$ линейному функционалу на $k$-формах $\varphi$ в $\mathbb{R} n$. Здесь $v$ - единичный касательный к $M k$-вектор. Линейный функционал $T_{M}$ в супественном характеризует $M$. Кроме того, определен (на этот раз и при неориентированном $M$ ) нелинейный функционал - вариобразие $V_{M} \varphi=\int_{M}|\langle\varphi, v\rangle| d F$. Интегральные нормы (массы) $\left\|T_{M}\right\|$ и $\left\|V_{M}\right\|$ совпадают с П., то есть с $F(M)$. Включение класса кусочно гладких подмногообразий пространства $\mathbb{R} n$ в более общие классы потоков и вариобразий играет в вариационном исчислении такую же роль, как обобщенные решения в теории уравнений в частных производных .



Среди всех потоков и вариобразий выделяют широкие классы целочисленных потоков и вариобразий. Последние сохраняют многие геометрич. свойства подмногообразий, Напр., целочисленный поток есть поток, допускающий представление $T=\sum_{i=1}^{\infty} n_{i} T_{M_{i}}$, где $n_{i}$ - целые числа, а $M_{i}$ суть $C^{1}$-гладкие подмногооб́разия, при условии, что конечна масса $(k-1)$-мерного потока $d T$, определяемого равенством $d T(\varphi)=T(d \varphi)$. Массы целочисленных потоков и вариобразий можно рассматривать как обобщения понятия П. поверхности. И здесь есть связь с П. по Лебегу. Пусть кусочно гладкие отображения $f_{i}: M \rightarrow \mathbb{R} n$ равномерно сходятся $f_{i} \rightarrow f$ и $L(M, f)=\lim _{i \rightarrow \infty} L\left(M, f_{i}\right)$. Тогда соответствующие вариобразия $V_{f_{i}}$ слабо сходятся к нек-рому целочисленному вариобразию $V_{f}$, причем $\left\|V_{f}\right\|=L(M, f)$. Тем самым каждому $f: M \rightarrow \mathbb{K}^{n}$ с $L(M, f)<\infty$ естественно сопоставляется вариобразие $V_{f}$ с массой $L(M, f)$; на языке потоков об әтом см. в [13].



Лum.: [1] Л е б е г А., Об измерении величин, пер. с франц., 2 изд., М., 1960; [2] Энциклопедия элементарной математики, т. 5, М.. 1ө66; [3] И л ь и н В. А., П о зн н к Э. Г., Основы математического анализа, 3 изд., ч. 1, М., 1971; [4] Л о п ш и ц А. М..

Вычисление площадей ориентированных фигур, М., 1956; [5] Вычисление площадей ориентированных фигур, М., 1956; [5]

$\Pi$ о г р е л о в А. В., Дифференциальная гепметрия, 5 изд., П о г о р е л о в А. В., Дифференциальная гепметрия, 5 изд.,

М., 1969; [6] Р а ш е в с к и й П. К., Риманова геометрия и тензорный анализ, 3 изд., M., 1967; [7] C e s a r i L., Surface area, Princeton, 1956; [8] F e d e r e r H., Geometric measure theory, B. - Hdlb. - N. Y., 1969; [9] F e d e r e r H., "Bull. Amer. Math. Soc.", 1952, v. 58, 스 3, p. $306-78 ;$ [10] Б y p a r o Ю. Д., 3 a л г а л л е p В. А., Геометрические неравенства, JI., 1980,

[11] B u s em a n $\mathrm{H}$ Н., "Ann. Math. 2 ser.", 1947, v. 48, p. 234-67; [12] A $1 \mathrm{mg} \mathrm{r}$ e $\mathrm{n}$ F. J., The theory of' varifolds, Princeton, 1965; [13] F e d e r e r H., "Trans. Amer. Math. Soc.", 1961, v. 98 , No 2 , p. $204-33$.



ПЛГККЕРА. Д.Бураго, В. А. Залгаллер, Л. Д. Кудрявчев. - модель, реализуюная геометрию трехмерного проективного пространства $P_{3}$ в гиперболлич. пространстве ${ }^{3} S_{5} \cdot \Pi$. и. основана на специальном истолковании плюккеровы $x$ координат прямой, определяемых для любой прямой пространства $P_{3}$.



При проективных преобразованиях пространства $P_{3}$ плюккөровы координаты преобразуются линейно. C помощью плюккеровых координат прямых про- странства $P_{3}$ устанавливается взаимно однозначное соответствие между прямыми пространствами $P_{3}$ и точками проективного пространства $P_{5}$, координаты к-рых численно равны плюккеровым координатам прямых пространства $P_{3}$.



Прямые пространства $P_{3}$ изображаются точками невырожденной квадрики пространства $P_{5}$, индекс к-рой равен трем.



Если считать эту квадрику за абсолют и определить в пространстве $P_{5}$ проективную (неевклидову) метрику, то получается пятимерное гиперболич. пространство ${ }^{3} S_{5}$. При каждой коллинеации и корреляции пространства $P_{3}$ происходит линейное преобразование плюккеровых координат, т. е. каждая коллинеация и корреляция изображаготся коллинеацией пространства $P_{6}$, переводящей в себя абсолют. Эти коллинеации тем самым являются перемещениями пространства ${ }^{3} S_{5}$. Перемещения пространства ${ }^{3} S_{5}$ изображают или коллинеации или корреляции пространства $P_{3}$.



Каждому линейному комплексу пространства $P_{3}$ сопоставляется точка пространства ${ }^{3} S_{5}$. Проективную геометрию пространства $P_{3}$ можно рассматривать как неевклидову геометрию гиперболич. пространства ${ }^{3} S_{5}$. Именно эта интерпретация геометрии пространства $P_{3}$ в пространстве ${ }^{3} S_{5}$ наз. и н т е р п р е т а ц и е й П л к к е р в связи с ролью плюккеровых координат.



Если в качестве основного образа пространства $P_{3}$ берется прямая, то геометрию әтого пространства можно рассматривать как геометрию на абсолюте пространства ${ }^{3} S_{5}$.



Группа проективных преобразований пространства $P_{3}$ изоморфна группе перемещений пространства ${ }^{3} S_{5}$, и всякому инволюционному проективному преобразованию пространства $P_{3}$ соответствует инволюционное перемещение пространства ${ }^{3} S_{5}$. Напр., пуль-системе в пространстве $P_{3}$ соответствует отражение от точки и от ее полярной гиперплоскости в пространстве ${ }^{3} S_{5}$ инволюционной гомологии в пространстве $P_{3}$ соот ветствует гиперболический паратактич. сдвиг на полупрямую в пространстве ${ }^{3} S_{5}$ и т. д. Каждой связной компоненте группы проективных преобразований пространства $P_{3}$ соответствует связная компонента группы перемещений пространства ${ }^{3} S_{5}$.



\begin{enumerate}

  \item Коллинеациям пространства $P_{3}$ с положительным определителем, включая тождественное преобразование, отвечают перемещения пространства ${ }^{3} S_{6}$ с определителем, равным +1 (сюда включаются тождественные преобразования).



  \item Корреляциям пространства $P_{3}$ с положительным определителем (включая нуль-систему) отвечают перемещения пространства ${ }^{3} S_{5}$ с определителем, равным - 1, переводящие собственную и идеальную области соответственно в себя (включая отражения от точки).



  \item Коллинеациям пространства $P_{3}$ с отрицательным определителем отвечают перемещения пространства ${ }^{3} S_{5}$ с определителем, равным +1 , переводящие собственную область в идеальную и обратно, эта компонента содержит гиперболич. сдвиг на полупрямую.



  \item Корреляциям пространства $P_{3}$ с отрицательным определителем отвечают перемещения пространства ${ }^{3} S_{5}$ с определителем, равным-1, переводящие собственную область в идеальную и обратно.



\end{enumerate}



Образам симметрии, соответствующим в пространствах $P_{3}$ и ${ }^{3} S_{5}$, сопоставляются числовые инварианты, между к-рыми существуют определенные связи.



II. и. применяется в исследованиях групп перемещений трехмерных неевклидовых пространств $S_{3}$, ${ }^{1} S_{3},{ }^{2} S_{3}$, к-рые изоморфны определенным подгруппам группы перемещений пространства ${ }^{3} S_{5}$. Устанавливается также взаимосвязь групп движений этих трехмерных пространств (эллиптического, гиперболиче-





\end{document}